%%% Local Variables:
%%% mode: latex
%%% TeX-master: "../main"
%%% End:

\chapter[结论与展望]{结论与展望}

\section{全文工作总结}

本文以两轮足机器人为研究对象，面向多地形、多步态、多算法的综合运动控制需求，提出了"分训-整合"两阶段的多任务强化学习框架，构建了由7个步态专家与1个学习式统一调度器组成的统一系统，在混合地形统一环境中实现了按地形自动切换步态，并系统性地研究了灾难性遗忘、奖励冲突与局部最优三大核心问题的应对策略。全文主要工作如下：

（1）建立了两轮足机器人的机构、运动学与倒立摆动力学模型，阐明了其欠驱动不稳定的控制本质；介绍了PPO、SAC、N3PO等强化学习算法与CMDP约束优化思想，并引入了分层强化学习与门控机制，为多步态控制与统一调度研究奠定理论与工具基础。

（2）提出了"分训-整合"两阶段的多任务强化学习总体框架，设计了包含平面、粗糙、碎石、斜坡、台阶、独木桥、木桩与狭窄通道等单元的混合地形统一仿真环境，明确了各典型地形与最适步态的对应关系，并从参数解耦、奖励分离、统一评估三个层面设计了三大核心问题的应对策略。

（3）完成了行走类步态专家的独立训练，实现了滚动行走（含侧向移动）与粗糙地形行走两类专家网络，覆盖稳定站立、腿长变化、前进后退、转向、侧向移动与地形自适应行走等基础能力。

（4）完成了精细平衡类步态专家的独立训练，实现了抬腿行走、抬腿上台阶与单腿平衡三类专家网络，其中抬腿上台阶采用N3PO约束强化学习在满足力矩与姿态约束的前提下实现稳定跨越。

（5）完成了高动态敏捷类步态专家的独立训练，构建了单PPO、PPO+SLIP、SLIP+PPO+VMC三种跳跃算法范式并系统对比了稳定性、收敛速度与追踪误差，提出了基于有限状态机与确定性翻转的后空翻训练方案。

（6）设计了学习式统一调度器，以地形感知与本体状态为输入，通过分层强化学习在混合地形统一环境中训练，结合滞回切换与能量指标辅助判别实现按地形自适应、平滑且稳定的步态切换。

（7）将7个步态专家与学习式统一调度器集成为完整系统，在混合地形统一环境中完成了整体遍历验证，并从灾难性遗忘、奖励冲突与局部最优三个维度开展了消融实验，验证了"分训-整合"范式的整体有效性与鲁棒性。

\section{主要创新点}

本文的主要创新点如下：

（1）\textbf{提出"分训-整合"两阶段多任务强化学习范式}。将综合任务分解为7类步态子任务分别独立训练步态专家，再通过学习式调度器统一整合，通过专家参数解耦从结构上规避了多步态连续学习中的灾难性遗忘，实现了两轮足机器人在混合地形环境中的完整遍历与按地形自动切换步态。

（2）\textbf{设计学习式统一调度器}。将"按地形选择步态"建模为分层强化学习问题，调度器以地形感知与本体状态为输入、以7个步态专家为选项集合，在混合地形统一环境中训练，并引入滞回切换机制与自适应能量指标辅助判别，实现了平滑、稳定、可泛化的步态自动切换。

（3）\textbf{构建专家奖励与调度器目标分离的分层奖励结构}。将各专家内部自洽的独立奖励函数与调度器高层目标解耦，结构性消解了"高动态-稳姿态"等子目标间的奖励冲突，并通过阶段式奖励塑形进一步细化高动态动作各阶段的目标聚焦。

（4）\textbf{构建多算法融合范式并解决高动态动作收敛难题}。系统对比单PPO、PPO+SLIP、SLIP+PPO+VMC三种跳跃算法范式，并针对后空翻提出基于有限状态机的前馈驱动确定性翻转方案，显著提升高动态动作的训练稳定性与可复现性。

（5）\textbf{引入N3PO约束强化学习实现精细平衡控制}。将抬腿上台阶建模为CMDP，通过N3PO在力矩与姿态约束下优化抬腿性能，结合收膝折腿与机身侧倾机制实现单腿平衡过独木桥，在强约束高动态场景下验证了约束强化学习的有效性。

\section{研究展望}

本文的研究工作仍存在一定的局限性，未来可在以下方向进一步深入：

（1）\textbf{仿真到现实的迁移}：本文成果目前主要基于仿真环境验证，后续将在两轮足机器人实物平台上开展sim-to-real迁移研究，评估领域随机化与调度策略在真实环境中的控制性能。

（2）\textbf{专家库的扩展与增量学习}：进一步扩展步态专家库（如新增其他敏捷动作），并研究专家增量加入时调度器的在线扩展与知识保留机制，提升系统的可持续学习能力。

（3）\textbf{高层任务规划与感知融合}：将视觉感知、地图与高层任务规划与调度器耦合，构建面向真实任务的"感知-规划-调度-控制"闭环系统，提升机器人在未知环境中的自主作业能力。

（4）\textbf{算法效率与泛化能力提升}：探索更高效的约束强化学习算法与更大规模的并行训练范式，进一步提升训练效率与系统在新地形上的泛化鲁棒性。
